\documentclass[11pt,a4paper]{article}
% Preámbulo común de las resoluciones de TPs de IA.
% Uso, desde temaN/resolucion-tpM.tex:  % Preámbulo común de las resoluciones de TPs de IA.
% Uso, desde temaN/resolucion-tpM.tex:  % Preámbulo común de las resoluciones de TPs de IA.
% Uso, desde temaN/resolucion-tpM.tex:  \input{../plantilla/preambulo}
\usepackage[T1]{fontenc}
\usepackage[utf8]{inputenc}
\usepackage[spanish,es-noshorthands,es-tabla]{babel}
\usepackage{lmodern}
\usepackage{microtype}
\usepackage[margin=2.3cm]{geometry}
\usepackage{xcolor}
\usepackage{booktabs}
\usepackage{tabularx}
\usepackage{array}
\usepackage{enumitem}
\usepackage{tcolorbox}
\usepackage{titlesec}
\usepackage{amsmath,amssymb}
\usepackage{graphicx}
\usepackage{caption}
\usepackage{tikz}
\usepackage[hang,flushmargin]{footmisc}
\usepackage[hidelinks]{hyperref}

\usetikzlibrary{babel,arrows.meta,positioning,calc,fit,decorations.pathreplacing}
\tcbuselibrary{skins,breakable}

\definecolor{azul}{HTML}{1F4E79}
\definecolor{gris}{HTML}{6B7280}
\definecolor{ambar}{HTML}{B45309}
\definecolor{fondoazul}{HTML}{F3F7FB}
\definecolor{fondogris}{HTML}{F4F4F5}
\definecolor{fondoambar}{HTML}{FFF7ED}

\titleformat{\section}{\Large\bfseries\color{azul}}{\thesection}{0.6em}{}
\titleformat{\subsection}{\large\bfseries}{\thesubsection}{0.6em}{}

% Enunciado de la consigna (gris).
\newtcolorbox{enunciado}[1]{
  breakable, enhanced, colback=fondogris, colframe=gris,
  boxrule=0.4pt, left=6pt, right=6pt, top=4pt, bottom=4pt,
  title={\bfseries #1}, coltitle=white, colbacktitle=gris, fonttitle=\small, fontupper=\small}
% Ampliación de la respuesta (azul). No es una tcolorbox: las notas al pie dentro de
% una tcolorbox que se parte entre páginas se superponen con el contenido.
\newenvironment{extras}{%
  \par\medskip\noindent
  {\setlength{\fboxsep}{3pt}\colorbox{azul}{\color{white}\bfseries\small\ EXTRAS\ }}%
  \par\nobreak\vspace{1pt}{\color{azul}\hrule height 0.6pt}\nobreak\vspace{3pt}}
 {\par\vspace{3pt}{\color{azul}\hrule height 0.6pt}\medskip}
% Advertencias, ambigüedades del enunciado, erratas (ámbar).
\newtcolorbox{aviso}[1]{
  breakable, enhanced, colback=fondoambar, colframe=ambar,
  boxrule=0.5pt, left=6pt, right=6pt, top=5pt, bottom=5pt,
  title={\bfseries #1}, coltitle=white, colbacktitle=ambar, fonttitle=\small}

\setlist[itemize]{itemsep=2pt, topsep=3pt, parsep=0pt}
\setlist[enumerate]{itemsep=2pt, topsep=3pt, parsep=0pt}
\captionsetup{font=small, labelfont=bf}
\setlength{\parindent}{0pt}
\setlength{\parskip}{4pt}

% Citas textuales (en el idioma original, sin traducir).
\newcommand{\cita}[1]{``\textit{#1}''}

% Fuentes, siempre en nota al pie.
%   \catedra{Episodio I}{Agentes en Inteligencia Artificial}
%   \aima{25.1}{973}     \pmg{12.1}{443}
\newcommand{\catedra}[2]{\footnote{Cátedra, #1, «#2».}}
\newcommand{\aima}[2]{\footnote{AIMA, §#1, p.~#2.}}
\newcommand{\pmg}[2]{\footnote{PMG, §#1, p.~#2.}}
% Para varias fuentes en una sola nota: \fuentes{AIMA, §2.1, p.~34; PMG, §12.1, p.~443}
\newcommand{\fuentes}[1]{\footnote{#1.}}

% Aclaración inicial de abreviaturas.
\newcommand{\notafuentes}{%
  {\small\color{gris} Las fuentes van en notas al pie. \textbf{Cátedra}: diapositivas de A.~J.~García (2026).
  \textbf{AIMA}: Russell \& Norvig, \textit{Artificial Intelligence: A Modern Approach}, 3.ª ed. (2010).
  \textbf{PMG}: Poole, Mackworth \& Goebel, \textit{Computational Intelligence: A Logical Approach} (1998).\par}}

% Íconos reutilizables (pics de TikZ). Uso: \pic[scale=0.8] at (x,y) {robot};
\tikzset{
  pics/robot/.style={code={
    \draw[thick, fill=azul!15] (-0.35,0.15) rectangle (0.35,0.65);            % cabeza
    \fill[azul] (-0.15,0.42) circle (0.06) (0.15,0.42) circle (0.06);         % ojos
    \draw[thick] (0,0.65) -- (0,0.8); \fill[ambar] (0,0.83) circle (0.05);    % antena
    \draw[thick, fill=azul!15] (-0.45,-0.45) rectangle (0.45,0.12);           % cuerpo
    \draw[thick] (-0.45,-0.1) -- (-0.65,0.1) (0.45,-0.1) -- (0.65,0.1);       % brazos
    \draw[thick, fill=gris!40] (-0.25,-0.55) circle (0.12) (0.25,-0.55) circle (0.12); % ruedas
  }},
  pics/compu/.style={code={
    \draw[thick, fill=azul!15, rounded corners=1pt] (-0.55,-0.25) rectangle (0.55,0.5); % pantalla
    \node[font=\tiny\ttfamily\bfseries, text=azul] at (0,0.12) {</>};
    \draw[thick] (0,-0.25) -- (0,-0.4); \draw[thick] (-0.25,-0.4) -- (0.25,-0.4);       % pie
  }},
  pics/persona/.style={code={
    \draw[thick] (0,0.35) circle (0.1);
    \draw[thick] (0,0.25) -- (0,-0.05) (-0.15,0.15) -- (0.15,0.15)
                 (0,-0.05) -- (-0.12,-0.3) (0,-0.05) -- (0.12,-0.3);
  }},
  pics/arbol/.style={code={
    \draw[thick, brown!60!black] (0,-0.3) -- (0,0);
    \fill[green!50!black!70] (0,0.15) circle (0.2);
  }},
  pics/documento/.style={code={
    \draw[thick, fill=white] (-0.2,-0.28) -- (-0.2,0.28) -- (0.1,0.28) -- (0.2,0.18) -- (0.2,-0.28) -- cycle;
    \foreach \y in {0.1,0,-0.1} \draw[gris] (-0.12,\y) -- (0.12,\y);
  }},
  pics/sobre/.style={code={
    \draw[thick, fill=white] (-0.25,-0.17) rectangle (0.25,0.17);
    \draw[thick] (-0.25,0.17) -- (0,-0.02) -- (0.25,0.17);
  }},
}

\usepackage[T1]{fontenc}
\usepackage[utf8]{inputenc}
\usepackage[spanish,es-noshorthands,es-tabla]{babel}
\usepackage{lmodern}
\usepackage{microtype}
\usepackage[margin=2.3cm]{geometry}
\usepackage{xcolor}
\usepackage{booktabs}
\usepackage{tabularx}
\usepackage{array}
\usepackage{enumitem}
\usepackage{tcolorbox}
\usepackage{titlesec}
\usepackage{amsmath,amssymb}
\usepackage{graphicx}
\usepackage{caption}
\usepackage{tikz}
\usepackage[hang,flushmargin]{footmisc}
\usepackage[hidelinks]{hyperref}

\usetikzlibrary{babel,arrows.meta,positioning,calc,fit,decorations.pathreplacing}
\tcbuselibrary{skins,breakable}

\definecolor{azul}{HTML}{1F4E79}
\definecolor{gris}{HTML}{6B7280}
\definecolor{ambar}{HTML}{B45309}
\definecolor{fondoazul}{HTML}{F3F7FB}
\definecolor{fondogris}{HTML}{F4F4F5}
\definecolor{fondoambar}{HTML}{FFF7ED}

\titleformat{\section}{\Large\bfseries\color{azul}}{\thesection}{0.6em}{}
\titleformat{\subsection}{\large\bfseries}{\thesubsection}{0.6em}{}

% Enunciado de la consigna (gris).
\newtcolorbox{enunciado}[1]{
  breakable, enhanced, colback=fondogris, colframe=gris,
  boxrule=0.4pt, left=6pt, right=6pt, top=4pt, bottom=4pt,
  title={\bfseries #1}, coltitle=white, colbacktitle=gris, fonttitle=\small, fontupper=\small}
% Ampliación de la respuesta (azul). No es una tcolorbox: las notas al pie dentro de
% una tcolorbox que se parte entre páginas se superponen con el contenido.
\newenvironment{extras}{%
  \par\medskip\noindent
  {\setlength{\fboxsep}{3pt}\colorbox{azul}{\color{white}\bfseries\small\ EXTRAS\ }}%
  \par\nobreak\vspace{1pt}{\color{azul}\hrule height 0.6pt}\nobreak\vspace{3pt}}
 {\par\vspace{3pt}{\color{azul}\hrule height 0.6pt}\medskip}
% Advertencias, ambigüedades del enunciado, erratas (ámbar).
\newtcolorbox{aviso}[1]{
  breakable, enhanced, colback=fondoambar, colframe=ambar,
  boxrule=0.5pt, left=6pt, right=6pt, top=5pt, bottom=5pt,
  title={\bfseries #1}, coltitle=white, colbacktitle=ambar, fonttitle=\small}

\setlist[itemize]{itemsep=2pt, topsep=3pt, parsep=0pt}
\setlist[enumerate]{itemsep=2pt, topsep=3pt, parsep=0pt}
\captionsetup{font=small, labelfont=bf}
\setlength{\parindent}{0pt}
\setlength{\parskip}{4pt}

% Citas textuales (en el idioma original, sin traducir).
\newcommand{\cita}[1]{``\textit{#1}''}

% Fuentes, siempre en nota al pie.
%   \catedra{Episodio I}{Agentes en Inteligencia Artificial}
%   \aima{25.1}{973}     \pmg{12.1}{443}
\newcommand{\catedra}[2]{\footnote{Cátedra, #1, «#2».}}
\newcommand{\aima}[2]{\footnote{AIMA, §#1, p.~#2.}}
\newcommand{\pmg}[2]{\footnote{PMG, §#1, p.~#2.}}
% Para varias fuentes en una sola nota: \fuentes{AIMA, §2.1, p.~34; PMG, §12.1, p.~443}
\newcommand{\fuentes}[1]{\footnote{#1.}}

% Aclaración inicial de abreviaturas.
\newcommand{\notafuentes}{%
  {\small\color{gris} Las fuentes van en notas al pie. \textbf{Cátedra}: diapositivas de A.~J.~García (2026).
  \textbf{AIMA}: Russell \& Norvig, \textit{Artificial Intelligence: A Modern Approach}, 3.ª ed. (2010).
  \textbf{PMG}: Poole, Mackworth \& Goebel, \textit{Computational Intelligence: A Logical Approach} (1998).\par}}

% Íconos reutilizables (pics de TikZ). Uso: \pic[scale=0.8] at (x,y) {robot};
\tikzset{
  pics/robot/.style={code={
    \draw[thick, fill=azul!15] (-0.35,0.15) rectangle (0.35,0.65);            % cabeza
    \fill[azul] (-0.15,0.42) circle (0.06) (0.15,0.42) circle (0.06);         % ojos
    \draw[thick] (0,0.65) -- (0,0.8); \fill[ambar] (0,0.83) circle (0.05);    % antena
    \draw[thick, fill=azul!15] (-0.45,-0.45) rectangle (0.45,0.12);           % cuerpo
    \draw[thick] (-0.45,-0.1) -- (-0.65,0.1) (0.45,-0.1) -- (0.65,0.1);       % brazos
    \draw[thick, fill=gris!40] (-0.25,-0.55) circle (0.12) (0.25,-0.55) circle (0.12); % ruedas
  }},
  pics/compu/.style={code={
    \draw[thick, fill=azul!15, rounded corners=1pt] (-0.55,-0.25) rectangle (0.55,0.5); % pantalla
    \node[font=\tiny\ttfamily\bfseries, text=azul] at (0,0.12) {</>};
    \draw[thick] (0,-0.25) -- (0,-0.4); \draw[thick] (-0.25,-0.4) -- (0.25,-0.4);       % pie
  }},
  pics/persona/.style={code={
    \draw[thick] (0,0.35) circle (0.1);
    \draw[thick] (0,0.25) -- (0,-0.05) (-0.15,0.15) -- (0.15,0.15)
                 (0,-0.05) -- (-0.12,-0.3) (0,-0.05) -- (0.12,-0.3);
  }},
  pics/arbol/.style={code={
    \draw[thick, brown!60!black] (0,-0.3) -- (0,0);
    \fill[green!50!black!70] (0,0.15) circle (0.2);
  }},
  pics/documento/.style={code={
    \draw[thick, fill=white] (-0.2,-0.28) -- (-0.2,0.28) -- (0.1,0.28) -- (0.2,0.18) -- (0.2,-0.28) -- cycle;
    \foreach \y in {0.1,0,-0.1} \draw[gris] (-0.12,\y) -- (0.12,\y);
  }},
  pics/sobre/.style={code={
    \draw[thick, fill=white] (-0.25,-0.17) rectangle (0.25,0.17);
    \draw[thick] (-0.25,0.17) -- (0,-0.02) -- (0.25,0.17);
  }},
}

\usepackage[T1]{fontenc}
\usepackage[utf8]{inputenc}
\usepackage[spanish,es-noshorthands,es-tabla]{babel}
\usepackage{lmodern}
\usepackage{microtype}
\usepackage[margin=2.3cm]{geometry}
\usepackage{xcolor}
\usepackage{booktabs}
\usepackage{tabularx}
\usepackage{array}
\usepackage{enumitem}
\usepackage{tcolorbox}
\usepackage{titlesec}
\usepackage{amsmath,amssymb}
\usepackage{graphicx}
\usepackage{caption}
\usepackage{tikz}
\usepackage[hang,flushmargin]{footmisc}
\usepackage[hidelinks]{hyperref}

\usetikzlibrary{babel,arrows.meta,positioning,calc,fit,decorations.pathreplacing}
\tcbuselibrary{skins,breakable}

\definecolor{azul}{HTML}{1F4E79}
\definecolor{gris}{HTML}{6B7280}
\definecolor{ambar}{HTML}{B45309}
\definecolor{fondoazul}{HTML}{F3F7FB}
\definecolor{fondogris}{HTML}{F4F4F5}
\definecolor{fondoambar}{HTML}{FFF7ED}

\titleformat{\section}{\Large\bfseries\color{azul}}{\thesection}{0.6em}{}
\titleformat{\subsection}{\large\bfseries}{\thesubsection}{0.6em}{}

% Enunciado de la consigna (gris).
\newtcolorbox{enunciado}[1]{
  breakable, enhanced, colback=fondogris, colframe=gris,
  boxrule=0.4pt, left=6pt, right=6pt, top=4pt, bottom=4pt,
  title={\bfseries #1}, coltitle=white, colbacktitle=gris, fonttitle=\small, fontupper=\small}
% Ampliación de la respuesta (azul). No es una tcolorbox: las notas al pie dentro de
% una tcolorbox que se parte entre páginas se superponen con el contenido.
\newenvironment{extras}{%
  \par\medskip\noindent
  {\setlength{\fboxsep}{3pt}\colorbox{azul}{\color{white}\bfseries\small\ EXTRAS\ }}%
  \par\nobreak\vspace{1pt}{\color{azul}\hrule height 0.6pt}\nobreak\vspace{3pt}}
 {\par\vspace{3pt}{\color{azul}\hrule height 0.6pt}\medskip}
% Advertencias, ambigüedades del enunciado, erratas (ámbar).
\newtcolorbox{aviso}[1]{
  breakable, enhanced, colback=fondoambar, colframe=ambar,
  boxrule=0.5pt, left=6pt, right=6pt, top=5pt, bottom=5pt,
  title={\bfseries #1}, coltitle=white, colbacktitle=ambar, fonttitle=\small}

\setlist[itemize]{itemsep=2pt, topsep=3pt, parsep=0pt}
\setlist[enumerate]{itemsep=2pt, topsep=3pt, parsep=0pt}
\captionsetup{font=small, labelfont=bf}
\setlength{\parindent}{0pt}
\setlength{\parskip}{4pt}

% Citas textuales (en el idioma original, sin traducir).
\newcommand{\cita}[1]{``\textit{#1}''}

% Fuentes, siempre en nota al pie.
%   \catedra{Episodio I}{Agentes en Inteligencia Artificial}
%   \aima{25.1}{973}     \pmg{12.1}{443}
\newcommand{\catedra}[2]{\footnote{Cátedra, #1, «#2».}}
\newcommand{\aima}[2]{\footnote{AIMA, §#1, p.~#2.}}
\newcommand{\pmg}[2]{\footnote{PMG, §#1, p.~#2.}}
% Para varias fuentes en una sola nota: \fuentes{AIMA, §2.1, p.~34; PMG, §12.1, p.~443}
\newcommand{\fuentes}[1]{\footnote{#1.}}

% Aclaración inicial de abreviaturas.
\newcommand{\notafuentes}{%
  {\small\color{gris} Las fuentes van en notas al pie. \textbf{Cátedra}: diapositivas de A.~J.~García (2026).
  \textbf{AIMA}: Russell \& Norvig, \textit{Artificial Intelligence: A Modern Approach}, 3.ª ed. (2010).
  \textbf{PMG}: Poole, Mackworth \& Goebel, \textit{Computational Intelligence: A Logical Approach} (1998).\par}}

% Íconos reutilizables (pics de TikZ). Uso: \pic[scale=0.8] at (x,y) {robot};
\tikzset{
  pics/robot/.style={code={
    \draw[thick, fill=azul!15] (-0.35,0.15) rectangle (0.35,0.65);            % cabeza
    \fill[azul] (-0.15,0.42) circle (0.06) (0.15,0.42) circle (0.06);         % ojos
    \draw[thick] (0,0.65) -- (0,0.8); \fill[ambar] (0,0.83) circle (0.05);    % antena
    \draw[thick, fill=azul!15] (-0.45,-0.45) rectangle (0.45,0.12);           % cuerpo
    \draw[thick] (-0.45,-0.1) -- (-0.65,0.1) (0.45,-0.1) -- (0.65,0.1);       % brazos
    \draw[thick, fill=gris!40] (-0.25,-0.55) circle (0.12) (0.25,-0.55) circle (0.12); % ruedas
  }},
  pics/compu/.style={code={
    \draw[thick, fill=azul!15, rounded corners=1pt] (-0.55,-0.25) rectangle (0.55,0.5); % pantalla
    \node[font=\tiny\ttfamily\bfseries, text=azul] at (0,0.12) {</>};
    \draw[thick] (0,-0.25) -- (0,-0.4); \draw[thick] (-0.25,-0.4) -- (0.25,-0.4);       % pie
  }},
  pics/persona/.style={code={
    \draw[thick] (0,0.35) circle (0.1);
    \draw[thick] (0,0.25) -- (0,-0.05) (-0.15,0.15) -- (0.15,0.15)
                 (0,-0.05) -- (-0.12,-0.3) (0,-0.05) -- (0.12,-0.3);
  }},
  pics/arbol/.style={code={
    \draw[thick, brown!60!black] (0,-0.3) -- (0,0);
    \fill[green!50!black!70] (0,0.15) circle (0.2);
  }},
  pics/documento/.style={code={
    \draw[thick, fill=white] (-0.2,-0.28) -- (-0.2,0.28) -- (0.1,0.28) -- (0.2,0.18) -- (0.2,-0.28) -- cycle;
    \foreach \y in {0.1,0,-0.1} \draw[gris] (-0.12,\y) -- (0.12,\y);
  }},
  pics/sobre/.style={code={
    \draw[thick, fill=white] (-0.25,-0.17) rectangle (0.25,0.17);
    \draw[thick] (-0.25,0.17) -- (0,-0.02) -- (0.25,0.17);
  }},
}


\title{\vspace{-1.5cm}\textbf{Trabajo Práctico 5 --- Resolución}\\[3pt]
\large Conceptos básicos de Robótica y Roboética}
\author{\normalsize Inteligencia Artificial --- UNS --- 2026}
\date{}

\begin{document}
\maketitle
\vspace{-1cm}
\notafuentes

%======================================================================
\section*{Ejercicio 1}

\begin{enunciado}{Consigna}
¿Qué es un robot y qué diferencias tiene frente a un agente de software? Brinde algún ejemplo que muestre claramente la diferencia.
\end{enunciado}

\paragraph{Robot.} Un agente físico que lleva a cabo tareas manipulando el mundo físico: percibe su entorno con \textbf{sensores} (cámara, sonar, GPS, acelerómetro) y actúa sobre él con \textbf{efectores} (ruedas, piernas, brazos, pinzas, hélices).\fuentes{Cátedra, Robótica y roboética, «Agentes físicos». Misma definición en AIMA, §25.1, p.~971}

\paragraph{Diferencia con un agente de software.} Los dos son agentes: entidades autónomas que perciben su entorno con sensores y actúan sobre él con efectores.\catedra{Episodio I}{Agentes en Inteligencia Artificial} Las diferencias están en el entorno:
\begin{itemize}
  \item \textbf{Qué entorno y con qué sensores y efectores}:\fuentes{Misma idea en AIMA, §2.1, p.~34}
  \begin{itemize}
    \item robot: el mundo físico, con sensores y efectores físicos;
    \item agente de software: un entorno de software (sistema operativo, archivos, Internet). Percibe teclas, archivos y paquetes de red, y actúa mostrando en pantalla, escribiendo archivos y enviando mensajes.
  \end{itemize}
  \item \textbf{No es una diferencia de complejidad}: un agente de software en Internet puede tener un entorno tan complejo como el mundo físico.\aima{2.3.1}{41}
  \item \textbf{Las limitaciones del mundo físico}: el robot real tiene que lidiar con un entorno parcialmente observable (la cámara no ve detrás de una esquina), estocástico (los motores patinan y tienen fricción), dinámico y continuo. Además, sus errores tienen consecuencias físicas: un choque real daña.\aima{25.1}{973}
\end{itemize}

\begin{center}
\small
\begin{tabularx}{\linewidth}{@{}>{\bfseries}l X X@{}}
\toprule
 & \textbf{Robot} & \textbf{Agente de software} \\
\midrule
Entorno & mundo físico & sistema operativo, archivos, Internet \\
Sensores & físicos: cámara, sonar, GPS & teclas, archivos, paquetes de red \\
Efectores & físicos: ruedas, brazos, pinzas & pantalla, escritura de archivos, mensajes \\
Complejidad & puede ser alta & puede ser alta (Internet) \\
Limitaciones físicas & parcialmente observable, estocástico, dinámico, continuo & --- \\
Si se equivoca & consecuencias físicas (un choque daña) & error en los datos \\
\bottomrule
\end{tabularx}
\end{center}

\begin{figure}[h]
\centering
\begin{tikzpicture}[
    caja/.style={draw=gris, thick, rounded corners=3pt, minimum width=3.6cm, minimum height=2.3cm},
    flecha/.style={-{Stealth[length=2.2mm]}, thick, draw=azul},
    rot/.style={font=\scriptsize, align=center, fill=white, inner sep=1pt}]
  % --- Robot ---
  \node[font=\bfseries\small, text=azul] at (4,2.1) {Robot};
  \node[caja, fill=green!5] (m1) at (0,0) {};
  \node[font=\scriptsize\bfseries, anchor=north] at (m1.north) {Mundo físico};
  \pic[scale=0.9] at (-0.8,-0.35) {arbol};
  \pic[scale=0.9] at (0,-0.35) {persona};
  \pic[scale=0.9] at (0.8,-0.35) {arbol};
  \node[caja, fill=fondoazul] (a1) at (8,0) {};
  \node[font=\scriptsize\bfseries, anchor=north] at (a1.north) {Agente};
  \pic[scale=0.9] at (8,-0.3) {robot};
  \draw[flecha] (m1.north east) ++(0,-0.4) -- node[rot, above] {sensores\\cámara, sonar, GPS} ($(a1.north west)+(0,-0.4)$);
  \draw[flecha] (a1.south west) ++(0,0.4) -- node[rot, below] {efectores\\ruedas, brazos, pinzas} ($(m1.south east)+(0,0.4)$);
  % --- Agente de software ---
  \node[font=\bfseries\small, text=azul] at (4,-2.6) {Agente de software};
  \node[caja, fill=fondogris] (m2) at (0,-4.7) {};
  \node[font=\scriptsize\bfseries, anchor=north] at (m2.north) {Entorno de software};
  \pic at (-0.8,-5) {documento};
  \pic at (0,-5) {sobre};
  \pic at (0.8,-5) {documento};
  \node[caja, fill=fondoazul] (a2) at (8,-4.7) {};
  \node[font=\scriptsize\bfseries, anchor=north] at (a2.north) {Agente};
  \pic at (8,-4.95) {compu};
  \draw[flecha] (m2.north east) ++(0,-0.4) -- node[rot, above] {sensores\\teclas, archivos, paquetes} ($(a2.north west)+(0,-0.4)$);
  \draw[flecha] (a2.south west) ++(0,0.4) -- node[rot, below] {efectores\\pantalla, archivos, mensajes} ($(m2.south east)+(0,0.4)$);
\end{tikzpicture}
\caption{Los dos son agentes que perciben con sensores y actúan con efectores; cambia el entorno.}
\end{figure}

\paragraph{Ejemplo.} La misma tarea, ``entregar un pedido'', hecha por dos agentes:
\begin{itemize}
  \item \textbf{Robot de delivery}: percibe la vereda con cámaras y GPS, y se mueve con ruedas. No ve al peatón que aparece detrás de un auto estacionado y, si frena tarde, puede atropellarlo.
  \item \textbf{Agente de software de la tienda online}: percibe la orden de compra (un mensaje), consulta el stock, genera la factura y envía un mail de confirmación. Todo lo que hace es leer y escribir datos: si se equivoca, el error se corrige en los datos.
\end{itemize}

\begin{extras}
\begin{itemize}
  \item \textbf{Cuatro aspectos de los robots contemporáneos}:\catedra{Robótica y roboética}{Una definición de robot}
  \begin{itemize}
    \item \textbf{movilidad}: para funcionar en entornos humanos o peligrosos para los humanos;
    \item \textbf{interactividad}: gracias a sensores y efectores;
    \item \textbf{comunicación}: interfaces de reconocimiento y síntesis de voz;
    \item \textbf{autonomía}: capacidad de ``razonar'' y tomar sus propias decisiones sin control externo directo.
  \end{itemize}
  \item \textbf{Un entorno físico puede ser simple}: un robot que inspecciona piezas en una cinta transportadora puede suponer que la iluminación es siempre la misma, que solo pasan piezas conocidas y que tiene dos acciones (aceptar o rechazar).\aima{2.3.1}{41}
  \item \textbf{El tiempo real}: \cita{the real world stubbornly refuses to operate faster than real time}. En simulación, un algoritmo de aprendizaje por refuerzo aprende en pocas horas de CPU a partir de millones de pruebas; con un robot real, esas pruebas llevarían años.\aima{25.1}{973}
  \item \textbf{Qué controla realmente el agente}: \cita{What an agent can control is what message (command) it sends to its actuators}. Lo que pase después depende del cuerpo y del mundo.\pmg{12.1}{444}
  \item \textbf{Robot como agente con propósito}: \cita{A robot is an artificial purposive agent}. Actúa para lograr ciertos resultados; sus preferencias suelen ser las de su diseñador, aunque no las conozca: con esta definición, hasta un termostato es un robot.\fuentes{PMG, §12.1, pp.~443--444}
  \item \textbf{Cuerpo y controlador}: \cita{A robot is made up of a body plus a controller. The controller receives percepts from the body and sends commands to the body.} El cuerpo convierte estímulos en percepciones (sensores) y comandos en acciones (actuadores).\pmg{12.2}{444}
\end{itemize}

  \begin{center}
  \begin{tikzpicture}[scale=1.15, transform shape,
      b/.style={draw=gris, thick, rounded corners=2pt, font=\scriptsize, minimum height=6mm, fill=white, inner sep=3pt},
      f/.style={-{Stealth[length=1.8mm]}, thick, draw=azul},
      l/.style={font=\tiny, fill=fondoazul, inner sep=1pt}]
    \node[b, fill=green!5] (env) at (0,0) {Entorno};
    \node[b] (sen) at (3,0.6) {Sensores};
    \node[b] (act) at (3,-0.6) {Actuadores};
    \node[b, fill=fondoazul!60!azul!20] (con) at (6,0) {Controlador};
    \node[draw=azul, dashed, rounded corners=3pt, fit=(sen)(act)(con), inner sep=5pt,
          label={[font=\tiny\bfseries, text=azul]above:Robot}] {};
    \node[draw=gris, densely dashed, rounded corners=2pt, fit=(sen)(act), inner sep=3pt,
          label={[font=\tiny, text=gris]left:cuerpo}] {};
    \draw[f] (env) |- node[l, pos=0.75] {estímulos} (sen);
    \draw[f] (sen) -| node[l, pos=0.25] {percepciones} (con);
    \draw[f] (con) |- node[l, pos=0.75] {comandos} (act);
    \draw[f] (act) -| node[l, pos=0.25] {acciones} (env);
  \end{tikzpicture}
  \end{center}
\end{extras}

\end{document}
